\chapter{System Architecture and Implementation Design}
\label{ch:architecture}

\section{End-to-end component architecture}

The repository uses versioned, append-only stages. Source BIDS and EEGLAB files
feed deterministic preprocessing partitions; tensor preparation creates new
arrays and indexes; each benchmark freezes its protocol, persists fold-level
models and predictions, and runs a separate verification/report step. Figure
\ref{fig:architecture} shows the implementation-level data flow.

\begin{figure}[!htbp]
\centering
\begin{tikzpicture}[node distance=9mm and 7mm,
block/.append style={text width=43mm,minimum height=17mm},
process/.append style={text width=43mm,minimum height=17mm}]
\node[block] (bids) {1. BIDS + EEGLAB\\\code{dataset/ds004504}};
\node[process,right=of bids] (prep) {2. Balancing and QC\\\code{preprocessing/}};
\node[block,right=of prep] (parts) {3. Versioned arrays\\primary, leftover, excluded\\and window indexes};
\node[process,below=of parts] (tensor) {4. Tensor preparation\\V1, V2, V3 arrays\\moments and controls};
\node[process,left=of tensor] (bench) {5. Nested benchmarks\\training-only bases\\and V4 tuning};
\node[block,left=of bench] (models) {6. Saved artifacts\\fold models, predictions\\and verification reports};
\draw[arrow] (bids)--(prep);
\draw[arrow] (prep)--(parts);
\draw[arrow] (parts)--(tensor);
\draw[arrow] (tensor)--(bench);
\draw[arrow] (bench)--(models);
\end{tikzpicture}
\caption{Implementation architecture. Follow steps 1--6 across the top row, then back along the bottom row. Derived stages use versioned directories; source files remain read-only.}
\label{fig:architecture}
\end{figure}


\section{Core preprocessing components}

\begin{description}[style=nextline]
  \item[\path{preprocessing/build_balanced_dataset.py}]
  Loads participant and recording metadata, down-samples derivatives, scans
  complete windows, constructs \FTD{} duration references, solves the SciPy MILP
  assignment, spreads chosen windows in time, writes signals/metadata/manifests,
  preserves leftovers and tails, and verifies the partition.

  \item[\path{preprocessing/audit_primary_signal_quality.py}]
  Computes robust per-window, per-channel, and per-subject signal statistics,
  spectral band summaries, shared-variance diagnostics, and candidate outlier
  flags without modifying the primary arrays.

  \item[\path{preprocessing/build_clean_primary_dataset.py}]
  Applies CAR and mean removal, assigns post-CAR exclusion reasons, performs
  deterministic medium-bucket balance trims, writes retained and excluded
  partitions, and verifies finiteness, group equality, means, and source
  partition completeness.
\end{description}

\section{Tensor and model components}

\begin{description}[style=nextline]
  \item[\path{tensor_pipeline/tensorize_subjects.py}]
  Builds the V1 non-negative relative-power tensor with 10 chronological rank
  bins and saves channel, frequency, subject, and label arrays.

  \item[\path{tensor_pipeline/run_decomposition_benchmark.py}]
  Implements PCA, NMF, non-negative CP, graph NCP, and non-negative Tucker,
  including held-out non-negative least-squares projection. The CPU launcher
  swaps the optional XGBoost classifier for histogram gradient boosting while
  retaining decomposition behavior.

  \item[\path{tensor_pipeline/v2/prepare.py and models.py}]
  Corrects electrode order, enforces boundary guards, constructs signed logPSD
  tensors and 418 controls, and implements reusable TT/HOSVD projections,
  anatomical smoothing, classifier grids, and trainable pipelines.

  \item[\path{tensor_pipeline/v2/benchmark.py}]
  Runs nested participant CV, fits decomposition bases using training people,
  selects configurations from inner out-of-fold predictions, persists bases and
  classifiers, and aggregates complete repeats.

  \item[\path{tensor_pipeline/v3_waveform/prepare.py}]
  Together with \path{wave_models.py}, this component
  rebuilds the full-cohort analysis index from primary plus leftovers, computes
  exact waveform-delay moments and same-cohort spectral controls, and implements
  unsupervised TT/Tucker energy projections, subject weights, and supervised TT.

  \item[\path{tensor_pipeline/v4_tuning/tune_models.py}]
  Creates the V4 feature bank, stronger-regularized TT objective, 19 classifier
  settings, seven representation settings, and six supervised TT settings.

  \item[\path{tensor_pipeline/v4_tuning/benchmark_tuning.py}]
  Freezes and evaluates 139 configurations in nested CV, produces automatic and
  per-family selections, stores training-only bases and full pipelines, and
  records every held-out prediction.
\end{description}

\section{Artifact schema}

Each major stage writes both numerical arrays and human/machine-readable
provenance:

\begin{itemize}
  \item \code{signals.npy} or feature arrays for efficient memory mapping;
  \item \code{metadata.npz} or \code{index.npz} for subject IDs, labels, sample
  IDs, window positions, counts, and channel names;
  \item JSON protocol, preparation, selection, verification, and fingerprint
  documents;
  \item fold-level JSON predictions and inner scores;
  \item trusted-local \code{joblib} models and bases; and
  \item portable XGBoost booster exports accompanied by their required complete
  preprocessing pipelines.
\end{itemize}

\section{Leakage barriers by construction}

The most important architectural boundary is participant membership. A
transformation $T$ is learned as $T_{\mathcal{I}_{\mathrm{train}}}$ from training
participants only, then applied unchanged to validation or test participants.
This rule covers decomposition means and bases, graph-pre-smoothed bases,
univariate feature selection, standardization, class weights, classifier fitting,
and supervised TT parameters.

In V2, exactly 24 windows per training subject are spread across that subject to
calibrate window-level decomposition bases. All of a held-out person's windows
are projected afterward. In V3/V4, per-subject delay moments and spectral
controls are fixed input summaries, while decomposition, feature selection, and
scaling are fit inside the relevant training fold.

\section{Computational design choices}

The completed runs are CPU-oriented. Memory maps prevent loading all waveform
windows simultaneously. Covariance eigensolves replace explicit huge unfolding
SVDs where mathematically equivalent. Delay moments compress waveform-filter
energy evaluation from all $19\times16\times485$ tensors to $304\times304$
matrices per subject. Thread pools are limited to four workers to avoid nested
parallel oversubscription. V4 uses XGBoost's CPU histogram tree method; the
installed build reports \code{USE\_CUDA=false}. PyTorch and the GPU were not
required for the completed experiments.
